\section{Discussion}\label{sec:discussion}

\subsection{What the results mean for the setting decision}\label{sec:disc-meaning}

The question a process engineer asks of a calibrated crusher model is how well it predicts a survey it was not fitted to. On the ten HP500 surveys
the answer has two parts. For throughput, the calibrated engine is a useful predictor: its held-out error of $12.09\%$ is close to its in-sample
error of $11.55\%$, it beats the constant mean by six points, and a nested $80\%$ band of about $\pm239$~t/h covers the measured feed rate at the
stated rate. The capacity physics, with one refit scalar, does as well as a free-form regression with two coefficients, which is what a physical
shape should buy on ten points. For power, the calibrated engine is not a useful predictor: the Bond draw it computes hardly varies across the
surveys while the measured power spans almost a factor of two, and the published power constants belong to a different net power. A plant that
wants to predict the power of this machine at a new setting is better served, on these data, by a line on the setting, which reaches $12.78\%$ held
out, than by any of the size-aware energy laws the product implements. The recommendation that follows is specific: use the calibrated engine for
throughput and for the product size trends it was built for, and read power from the measured relation with the setting until the plant measures
the quantities that the energy laws need (the crushing work index from an SMC Test, the no-load power, and the product size distribution of each
survey).

\subsection{Why each method behaves as it does}\label{sec:disc-why}

The setting-only line wins on power because the measured power on these surveys is, to first order, a function of the setting ($r=-0.82$), and
the line spends its two coefficients on exactly that relation. The free-form fit spends a third coefficient on $f_{80}$, which adds noise on ten
points and costs it $1.25$ points held out. The Bond draw is flat because, as the setting tightens, the finer product raises the specific energy
while the lower throughput lowers the tonnage, and the two nearly cancel in kilowatts; the measured power instead rises as the setting tightens,
which a model of the machine's mechanics might explain, for example an energy computed as the product of the liner pressure and
the displacement over the chamber~\cite{liu2018}, which the product does not implement. Morrell's energy inherits the same design column, the specific energy times the measured throughput,
so its fit has the same problem as Bond's, and the variant with a no-load term recovers accuracy only by fitting a negative work index. The bounded
residual fails held out because it is a two-feature correction fitted on nine points to residuals that are themselves noisy: its leaky error is the
lowest of all predictors and its strict error is above the backbone's, the signature of overfitting. CrusherBayes calibrates its interval because
its noise and discrepancy scales are learned from the residuals of the engine (the posterior noise is about two and a half times the prior median on
both targets), and narrows it relative to the backbone band on power because its discrepancy slope on the setting absorbs the relation that the
Bond draw misses, about $57$~kW per standard deviation of the setting. That slope plays the role of the setting-only line inside the Bayesian model, which is
consistent with CrusherBayes approaching but not passing the line on power.

\subsection{Where each method fails}\label{sec:disc-failures}

Each method has a failure that the results expose. The Bond draw and the published power form fail on the measured power (Table~\ref{tab:power-parity}).
Morrell through the origin fails by $25$ points against the line, and with a no-load term it fits a physically impossible work index in every fold
(Table~\ref{tab:morrell}). The bounded residual fails to beat the backbone held out, and the label-shuffle control shows how a leaky evaluation would
have hidden that (Table~\ref{tab:leaky}). CrusherBayes does not beat the setting-only line on power, $14.31\%$ against $12.78\%$, and does not beat the
backbone on throughput beyond noise ($0.09$ points); its posterior leaves the classification and breakage scalars at their prior, so it cannot say
anything about the product size distribution. The calibrated backbone misses two surveys of ten on throughput, survey 1 by $407$~t/h, so a single
survey can sit well outside a model whose average error is twelve percent. The learned surrogate is worst on the fine tail of the product
($9.21\%$ on $P_{20}$), and the autoencoder misses $72\%$ of the planted off-manifold points. The circuit cases have no measured data, so their absolute
numbers carry no validation at all. The survey HP500-03 is flagged pass-through because its feed is barely coarser than its setting; a model reading
that survey's product size is outside its regime.

\subsection{What the product recommends and when}\label{sec:disc-recommend}

Three recommendations follow, each with the condition under which it holds. First, read the calibrated engine's throughput with its nested band
(about $\pm21\%$) for a secondary HP500 on this ore inside the surveyed support of $38$ to $55$~mm and $47$ to $140$~mm; outside it the envelope check
flags extrapolation and the band does not apply. Second, for power, prefer the measured relation with the setting, and use the CrusherBayes interval,
which covers the measured power at the stated rate and is narrower than the backbone band, when an uncertainty is needed; do not read the engine's
Bond power as a prediction of a survey. Third, to calibrate the population balance itself, which ten surveys without product sizes cannot do,
measure the product size distribution of each survey: the contraction table shows that no other addition to the present data would identify the
classification and breakage scalars. For the synthetic machine classes the product makes no recommendation beyond the trends; their constants are
typical of their class and their band on the secondary cone is a lower bound.

\section{Limitations}\label{sec:limits}

\emph{Sample size.} The real data are ten surveys of one machine on one ore. A difference below about three MAPE points is within noise, and one survey
moves a coverage by $0.1$, so the coverages of $0.80$ and $0.90$ are coarse estimates, and the ranking of predictors within three points of each other is
not stable. The surveys are treated as exchangeable; their order in time is not known to the product.

\emph{What is not measured.} The product size distribution of the surveys is printed only as figures in the source, so the product size is never a target,
and every statement about the engine's $P_{80}$ on the survey lane is a model output. The throw and the speed of each survey are not in the table and are
held at nominal values. The crushing work index and the no-load power of the machine are not published, so Morrell's method and the power form are fits.

\emph{The source calibration.} The constants of the calibrated lane are carried from the product's code and earlier documentation; the publisher page of
the source refused automated access in every documenting attempt, so their correspondence to the numbered tables of the paper, and the paper's own error
figures, are not re-verified here. The priors of CrusherBayes have widths chosen by the author, because no published range of these fitted scalars for this
ore was found.

\emph{Scope of the engine.} The engine is a single-zone population balance with a reduced capacity model and an empirical power law. It has no chamber
geometry, no liner wear, no feed segregation and no choke-level dynamics. The synthetic machines are not calibrated: their constants are illustrative
of their class and reproduce trends. The three-dimensional chamber view is a kinematic view of the geometry and not a particle simulation. The learned
models emulate the engine and not a plant, and the autoencoder detects distance from the engine's sampled envelope and not a mechanical fault. The drift
projection models no wear and asserts no drift rate.

\emph{Discrete-element simulation.} A discrete-element simulation with breakage, by particle replacement with the Tavares breakage model~\cite{tavares2022}
in a separate package, is being built for validation on the laboratory jaw crusher of Barrios et al.~\cite{barrios2020}; it is not part of the methods or the
results of this report. The open HP500 data could not validate it in any case, because the liner geometry of that machine is not public.

\emph{Prior work.} The searches that preceded CrusherBayes found no Bayesian calibration with structural discrepancy of a cone-crusher population-balance
model evaluated held out; they were not exhaustive (two citation indexes, Scopus and Web of Science, were not searched). Cone-crusher models that couple
energy and breakage~\cite{lindqvistli2021,bengtsson2019} and machine-learning models of cone-crusher power~\cite{gharaei2024} exist and were read only at the
level of their metadata; the product implements neither.
